\section{Evidence}\label{sec:evidence}
We first follow one data set through both groupings (\cref{sec:worked}), then report the large benchmark on random
equations and the development benchmark (\cref{sec:bench}). An answer is \emph{right} if the recovered set of factors
matches the true one within relative coefficient error $0.05$; \emph{equivalent} if it explains the data as well as the
true equation (a $\chi^2$ test adjusted for degrees of freedom, or the same product written differently); an
\emph{abstention} if the label is \IDK{} or \UND; and \emph{false-confident} if a wrong equation is returned as if it
were right. Library-regression baselines cannot abstain, so each of their wrong answers is false-confident.

\subsection{One run, step by step}\label{sec:worked}
The data are the running example with $c=1$, $\nu=0.3$ on a $64\times128$ grid over $[0,2]\times[0,2\pi)$, with white
noise of standard deviation $0.05\,\mathrm{std}\abs{u}$. \emph{Modes}: only $\kappa\le4$ rise above the detection
threshold; the matrix pencil returns twelve modes. \emph{Grouping I} (\cref{fig:directions}): eight modes lie on the line
$\operatorname{Im}\lambda=-\kappa$ with $\operatorname{Re}\lambda\approx0$, four on the parabola $\operatorname{Re}\lambda=-0.3\kappa^2$;
the transport modes occupy eight distinct wavenumbers and the diffusion modes four, against two and three required by
\cref{prop:counting}. \emph{Class}: the structure $((1,1),(1,2))$ ranks first. \emph{Grouping II}
(\cref{fig:grouping}): over $32$ posterior samples the $64$ factors form two tight clusters, with masses $0.767$ and
$0.233$ and $P(k=2)=1$. \emph{Components} (\cref{fig:components}): the transport component consists of tilted stripes, a
profile moving at speed $1$; the diffusion component has vertical structure fading in time; the residual, $2.5\%$ of the
data norm, shows no structure. \Cref{tab:worked} lists the output. The groups are only moderately separated
($\thetaF=0.60$, amplification $\gN^{-1/2}=2.41$), and the intervals still cover the truth; the two conjugate factors of
the beam equation, much closer to each other ($\cF=0.939$), are recovered with the same accuracy and an amplification of
$4.05$ (\SMcases).

\begin{figure}[tbp]
\centering
\includegraphics[width=\linewidth]{fig_components}
\caption{The worked example as functions of $(t,x)$ on a common colour scale: (a) noisy data; (b) transport component;
(c) diffusion component; (d) residual, magnified five times, indistinguishable from the noise.}\label{fig:components}
\end{figure}

\begin{table}[tbp]
\centering\small
\caption{Output of the worked example and how to read it. Intervals are $90\%$ intervals from $16$ parametric-bootstrap
refits.}\label{tab:worked}
\begin{tabularx}{\linewidth}{@{}l l X@{}}
\toprule
Quantity & Value & Meaning\\
\midrule
equation & $(\tau+1.0006\,\xi)(\tau-0.2980\,\xi^2)$ & recovered factors; truth $(\tau+\xi)(\tau-0.3\,\xi^2)$\\
coefficient of $\xi$ & $1.0006\ [0.9998,1.0020]$ & speed $c$ of the transport group\\
coefficient of $\xi^2$ & $-0.2980\ [-0.3049,-0.2919]$ & diffusivity $\nu$ of the diffusion group\\
label & \FC & enough distinct wavenumbers per factor; goodness of fit passed\\
$P(k)$; $P(\text{split})$ & $P(k=2)=1$; $1.000$ & number of groups; evidence that both groups are needed\\
$w_1,\ w_2$ & $0.768,\ 0.232$ & energy shares of the two groups\\
$\gN,\ \thetaF$; $\gN^{-1/2}$ & $0.173,\ 0.60$; $2.41$ & separation of the groups; worst-case error amplification\\
\bottomrule
\end{tabularx}
\end{table}

\subsection{Benchmarks}\label{sec:bench}
\paragraph{Large benchmark on random equations} To test generalization we drew $300$ equations at random from five
families, $60$ each: first-order factors (F1); first-order dispersive factors (F2); irreducible
second-order factors (F3); products of two or three factors (F4); and third-order factors or a reaction factor times a
transport factor, often with few excited wavenumbers (F5). Every draw is a different equation with an independent field
and independent noise ($1\%$ of $\mathrm{std}\abs{u}$, regular grid), and none was used to set any parameter: the pipeline
ran with the defaults fixed during development. We compare with weak SINDy~\cite{MesBor21}, the strongest library
regression on the development benchmark below, run through PySINDy with a linear library of mixed derivatives and a
sparsity threshold chosen by BIC. $285$ draws finished; $15$ (twelve in F4, three in F5) exceeded the time limit of three
hours per draw and are reported separately. \Cref{tab:large} and \cref{fig:evidence}(a) give the outcomes.

\begin{table}[tbp]
\centering\small
\caption{Large benchmark: $285$ random equations not used in development, grid, $1\%$ noise. ``Right+eq.'' is the share
of right or equivalent answers with a Wilson $95\%$ interval; ``abstain'' includes one correct abstention (F2) and one
crash (F4).}\label{tab:large}
\begin{tabular}{@{}llrrrrr@{\hspace{14pt}}l@{}}
\toprule
family & method & $n$ & right & equiv. & abstain & false-conf. & right+eq.\\
\midrule
all & this method & 285 & 187 & 88 & 10 & \textbf{0} & \textbf{96.5\%} [93.7, 98.1]\\
    & weak SINDy  & 285 &  77 &  4 &  0 & 204 & 28.4\% [23.5, 33.9]\\
\midrule
F1 & this method & 60 & 59 &  1 & 0 & 0 & 100\%\\
   & weak SINDy  & 60 & 42 &  0 & 0 & 18 & 70.0\%\\
F2 & this method & 60 & 31 & 23 & 6 & 0 & 90.0\%\\
   & weak SINDy  & 60 &  9 &  0 & 0 & 51 & 15.0\%\\
F3 & this method & 60 & 41 & 18 & 1 & 0 & 98.3\%\\
   & weak SINDy  & 60 & 26 &  2 & 0 & 32 & 46.7\%\\
F4 & this method & 48 & 12 & 34 & 2 & 0 & 95.8\%\\
   & weak SINDy  & 48 &  0 &  1 & 0 & 47 & 2.1\%\\
F5 & this method & 57 & 44 & 12 & 1 & 0 & 98.2\%\\
   & weak SINDy  & 57 &  0 &  1 & 0 & 56 & 1.8\%\\
\bottomrule
\end{tabular}
\end{table}

\begin{figure}[tbp]
\centering
\includegraphics[width=\linewidth]{fig_evidence}
\caption{Outcomes. (a) Large benchmark, overall and per family: share of draws that are right, equivalent, abstained or
false-confident, with the right-or-equivalent share at the end of each bar. (b) Development benchmark: $10$ equations
$\times$ $3$ designs. Hatched: a wrong equation returned without any warning.}\label{fig:evidence}
\end{figure}

The method is right or equivalent on $275$ of $285$ draws ($96.5\%$) and false-confident on none; the exact one-sided
$95\%$ upper bound on its false-confident rate is $1.0\%$. Weak SINDy is right or equivalent on $81$ ($28.4\%$) and
false-confident on $204$. On the same draws the method succeeds where weak SINDy fails $194$ times and the reverse never
happens (exact sign test, $p=8\cdot10^{-59}$). The gap is largest exactly where the grouping principle predicts it:
on products of factors (F4) and on higher-order and reaction factors (F5) weak SINDy is right on one draw in fifty,
the method on $96$--$98\%$. Where both are right, the method's median relative coefficient error is $0.4\%$, against
$1.4\%$. Counting every one of the $15$ unfinished draws as a failure, the method is still right or equivalent on
$275/300=91.7\%$ (Wilson interval $88.0$--$94.3\%$); the unfinished draws are the slowest, not a random subset.

\paragraph{Development benchmark} Ten classical equations (transport, heat, advection--diffusion, Airy, Schr\"odinger,
beam, wave, Klein--Gordon, telegraph and the running example), each in three designs: a $64\times128$ grid with noise $1\%$,
the same grid with $5\%$, and $400$ scattered points with $5\%$. The defaults were tuned on these equations, so this is not
a held-out test; it compares the method with five baselines on classical cases (\cref{fig:evidence}(b),
\cref{tab:bench}). The method recovers all $30$ factor sets. The best baseline, weak SINDy, is right or equivalent on $25$
and false-confident on $5$, four of them products of factors. Repeating the two grid designs with three new fields and
noise draws each ($60$ independent inputs), the method is right on all $60$ with coefficient coverage $0.94$ at nominal
$0.90$; weak SINDy is false-confident on $13$.

\begin{table}[tbp]
\centering\small
\caption{Development benchmark: right or equivalent / false-confident, out of $10$ inputs per design, and the
repetition with independent fields and noise (out of $60$). Runtimes are wall-clock medians per input.}\label{tab:bench}
\begin{tabular}{@{}lcccccc@{}}
\toprule
Method & grid 1\% & grid 5\% & 400 pts 5\% & total & independent & median time\\
\midrule
this method        & 10 / 0 & 10 / 0 & 10 / 0 & 30 / 0 & 60 / 0 & 127 s\\
weak SINDy         & 10 / 0 & \phantom{1}7 / 3 & \phantom{1}8 / 2 & 25 / 5 & 47 / 13 & 3 s\\
ensemble SINDy     & \phantom{1}9 / 1 & \phantom{1}7 / 3 & \phantom{1}6 / 4 & 22 / 8 & -- & 5 s\\
SBR SINDy          & \phantom{1}8 / 2 & \phantom{1}6 / 4 & \phantom{1}6 / 4 & 20 / 10 & -- & 14 s\\
PDE-FIND           & \phantom{1}1 / 9 & \phantom{1}1 / 9 & \phantom{1}3 / 7 & \phantom{1}5 / 25 & -- & 0.6 s\\
BOP-DMD            & \phantom{1}1 / 9 & \phantom{1}0 / 10 & \phantom{1}0 / 10 & \phantom{1}1 / 29 & -- & 1.5 s\\
\bottomrule
\end{tabular}
\end{table}

\paragraph{Parameters and cost} The user-facing parameters are few: the starting and largest class ($C(3,4)$ and
$C(4,8)$), the detection threshold $z=3$, the goodness-of-fit level $10^{-3}$, the evidence margin of one nat, and the
number of posterior samples and bootstrap refits ($8$ and $16$); \SMalgo{} lists all of them. All were fixed during
development and left unchanged for the large benchmark. The price of exploring a whole class of factored structures on
every sample is time: a median of $127$\,s per input on the development benchmark and $13$ minutes on the random
equations, against seconds for library regression.

\begin{table}[b]
\centering\small
\caption{The kinds of claim made in the paper.}\label{tab:status}
\begin{tabularx}{\linewidth}{@{}>{\raggedright\arraybackslash}p{2.0cm}X@{}}
\toprule
Kind & Statements\\
\midrule
proved & groups are factors with overlap $\deltaP$ (\cref{thm:factor}); terms are not groups (\cref{prop:terms});
stability via $\gN$ (\cref{thm:hm}) and its Bayesian form (\cref{prop:bayes-angle}); identification and its limits on
finite data (\cref{thm:ident,prop:certify}); counting condition (\cref{prop:counting})\\
design rule & cluster modes onto curves, not terms onto a library; labels instead of ``identified''; decisions on the raw
observations; growth of the class instead of coefficient penalties; minimum-norm split of the overlap\\
measured & $275/285$ right or equivalent and $0$ confident wrong answers on random equations not used in development;
$30/30$ on the development benchmark and $60/60$ on its independent repetitions; coefficient coverage $0.94$ at nominal
$0.90$\\
\bottomrule
\end{tabularx}
\end{table}

\paragraph{What these numbers do and do not show} All inputs are synthetic and the noise is white; the large benchmark
uses one noise level, gridded data and one baseline. The claims we draw are therefore qualitative but sharp: on
equations whose symbol is a product of factors, or has factors of higher order, grouping by factors recovers the
structure where regression on a library of terms does not, and when the data do not suffice the labels and intervals
warn instead of returning a wrong equation.

